\section{Finite-instance cyclic audit}
\label{sec:numerical-reconstruction}

Table~\ref{tab:representative-endpoints} separates the two reductions at the
problem endpoint.  The grid pays only a symmetry cost.  The random row has
trivial $G_B$ but pays a strictly cyclic cost: its symmetric ground state is
absent from $\mathcal K_u$.  The final row has neither cost, yet its cyclic
excitation gap is hundreds of times larger than the full gap.  Thus a
dimension deficit, a ground-energy cost, and a spectral-gap change are
logically distinct.
All finite-instance Hamiltonians reported here and in
Appendix~\ref{app:numerical-protocol} use the penalty energies $\lambda=5$
and $\mu=10$ in
Eq.~\eqref{eq:hprob-rewrite}.

At the problem endpoint $H_{\rm prob}$, write $\Delta_{\rm full}$,
$\Delta_{\rm sym}$, and $\Delta_{\rm cyc}$ for the excitation gaps on
$\mathcal H_{\rm addr}$, $\mathcal H_{\rm sym}$, and $\mathcal K_u$.
The energy costs $\delta_{\rm sym}$ and $\delta_{\rm cyc}$ are those in
Eq.~\eqref{eq:accessibility-cost-rewrite}.  If $\mathcal G_{\rm cyc}$ is
the cyclic ground eigenspace, define its worst-case cover probability by
\begin{equation}
 p_{\rm opt}^{\min}
 =\min_{\substack{\psi\in\mathcal G_{\rm cyc}\\\|\psi\|=1}}
 \langle\psi|P|\psi\rangle,
\end{equation}
where $P$ is the feasible-cover projector from
Section~\ref{sec:rewrite-localization}.

\begin{table*}[t]
\centering
\caption{Representative problem-endpoint reconstructions.  Dimensions and
spectra are computed in the physical full Hamiltonian and exact cyclic
restriction.  The complete eleven-instance audit is in
Appendix~\ref{app:numerical-protocol}.}
\label{tab:representative-endpoints}
\small
\begin{tabular}{lrrrrrrrr}
\hline
instance & $\dim\mathcal H_{\rm sym}$ & $\dim\mathcal K_u$ &
$\delta_{\rm sym}$ & $\delta_{\rm cyc}$ & $\Delta_{\rm full}$ &
$\Delta_{\rm sym}$ & $\Delta_{\rm cyc}$ & $p_{\rm opt}^{\min}$\\
\hline
grid-2x4 & 22 & 22 & 0.271 & 0 & 0.000382 & 0.315 & 0.315 & 0.9944\\
rnd-b8-t7-p50-s002 & 120 & 101 & 0 & 0.0581 & 0.0563 & 0.0581 & 0.128 & 0.9986\\
scpe1-b8-t30 & 792 & 792 & 0 & 0 & 0.00591 & 4.87 & 4.87 & 0.9949\\
\hline
\end{tabular}
\end{table*}

Across all eleven frozen instances, five have
$\mathcal K_u\subsetneq\mathcal H_{\rm sym}$, but only one has
$\delta_{\rm cyc}(1)>0$; eight have $\delta_{\rm sym}(1)>0$ and nine have
$\delta_{\rm acc}(1)>0$.  On the audited path grid,
$|\delta_{\rm cyc}(s)|\leq1.2\times10^{-14}$ for every $s\leq s_c$, as
predicted by Theorem~\ref{thm:shell-cyclic-closure}.  Every cyclic endpoint
ground state is rank one with $p_{\rm opt}^{\min}\geq0.9944$.  These are
finite diagnostics, not scaling evidence.  Appendix
\ref{app:numerical-protocol} gives the full tables, shell-closure
conditioning, residual-based degeneracy rule, irrep classification, and
OR-Library provenance.

\section{Discussion and conclusion}
\label{sec:rewrite-conclusion}

For a fixed initial state and a symmetry-preserving interpolation, the
physically relevant spectrum is neither the full spectrum nor, in general,
the entire group-fixed spectrum.  It is the spectrum of the Hamiltonian
restricted to the cyclic space generated by that protocol.  Thus
$\mathcal K_u\subseteq\mathcal H_{\rm sym}$ is the central distinction:
changing the initial state, the generated algebra, or a preserved symmetry
can change the accessible spectrum and therefore the applicable gap.  The
resulting inaccessibility statement is correspondingly conditional on state
preparation and on every symmetry preserved along the path.

Within this framework, orbit invariance identifies unlabelled candidate
covers as the natural sites of the fixed-sector quotient, while the shell
closure determines which amplitude patterns among those sites the endpoint
algebra can generate.  The sector-resolved localization proposition
certifies cover-measurement probability with the required kinetic floor.
Exact orbit quotients separate ground multiplicity, excitation gaps,
isolation gaps, and accessibility costs, and expose coincidences between
irreducible sectors that are dark to the symmetric protocol.  At the
hopping-cancellation point, the original interpolation has an exact global
multiplicity closure for every feasible fixed-cardinality instance.  The
even cycle is a cover-transitive illustration whose fixed excitation gap is
at least $\lambda/2$ at that point.  A distinct Johnson--Metropolis parent
path has a uniformly polynomial cyclic
gap bounded below by $1024\,n^{-13}$ along the path, and consequently a
conditional polynomial adiabatic runtime in the abstract
Hamiltonian-access model.

These results establish a structural separation, not a practical quantum
speedup.  The cycle instances are classically solvable, Dicke-state
preparation has a cost that must be counted, and the Gibbs parent endpoint is
not the original problem Hamiltonian.  Likewise, the address count
$k\lceil\log_2n_B\rceil$ is not a complete hardware resource estimate:
incidence access, flags, ancillas, and gate depth require a separate
implementation analysis.

A natural next step is a decorated family whose classical cover structure
is nontrivial while its joint orbit quotient, dark global branch, and
polynomial accessible gap remain stable under controlled
symmetry-preserving perturbations.  More broadly, spectral claims for
symmetric Hamiltonian algorithms should identify not only a Hamiltonian and
its global gap, but also the initial state, preserved symmetry, generated
cyclic space, and isolated band actually followed by the protocol.
